% !TEX root = ../main.tex
%----------------------------------------------------------
% 08-discussion.tex — Section 7: Discussion and Future Work
%----------------------------------------------------------

\section{Discussion and Future Work}\label{sec:discussion}

This paper has developed a self-contained verification framework for
possibilistic forecasts.  The core contribution is a five-number
scorecard (Section~\ref{sec:native_scorecard}) that evaluates a
subnormal possibility distribution on its own terms---without first
converting it to a probability.  Supporting this scorecard are a
possibility-to-probability bridge
(Section~\ref{sec:poss_prob_bridge}) that preserves the ignorance
signal as an explicit outcome rather than erasing it through
normalization, and three complementary verification lanes
(Section~\ref{sec:three_lanes}) that together cover
categorical, probabilistic, and native possibilistic evaluation.
Worked examples in Section~\ref{sec:worked_examples} demonstrated
that each uncertainty component---possibility, ignorance, and
conditional necessity---independently contributes verification value,
so that collapsing the distribution into a single point forecast
discards actionable information.  Although the SPC convective outlook
categories served as the running example throughout, the framework
applies to any finite universe of discourse $\Omega$ over which a
subnormal possibility distribution is defined.

%----------------------------------------------------------
\subsection{Connections to Dempster--Shafer Theory}%
\label{sec:ds_connections}

The possibility--necessity duality has a well-known structural
counterpart in Dempster--Shafer (DS) theory
\citep{Shafer1976-MISSING}.  In DS terms, the possibility measure
$\poss(A)$ corresponds to the plausibility function $\mathrm{Pl}(A)$:
both represent an upper bound on the probability of the event $A$.
Conditional necessity $\Nc(A)$ plays the role of the belief function
$\mathrm{Bel}(A)$: both represent a lower bound.  The interval
$[\Nc(A),\,\poss(A)]$ used throughout this paper is therefore
structurally equivalent to the DS interval
$[\mathrm{Bel}(A),\,\mathrm{Pl}(A)]$, and the verification metrics
defined in Section~\ref{sec:native_verification}
can be read as special cases of
interval-probability verification.

A key difference concerns the nature of the residual uncertainty.
Possibilistic ignorance $\ign = 1 - \max(\pi)$ quantifies a
\emph{coverage gap}: the fraction of plausibility that the forecasting
system left unassigned to any outcome.  In DS theory, uncertainty
instead arises from \emph{evidential conflict} among multiple
information sources, and is distributed across focal elements of
varying specificity.  The scalar $\ign$ therefore has no direct DS
analogue; it is a distinctly possibilistic quantity.  The
possibility-to-probability bridge
(Section~\ref{sec:poss_prob_bridge}) is inspired by the pignistic
transformation of \citet{Smets1990-MISSING} but extends it to
subnormal distributions: the classical pignistic transform assumes a
normal (i.e., $\max(\pi)=1$) mass function, whereas the bridge
explicitly allocates the subnormality residual $\ign$ to a dedicated
ignorance outcome before redistributing the remaining mass.

The structural parallels suggest a natural avenue for generalization.
The verification framework developed here could be extended to
belief-function forecasts by replacing $\poss$ and $\Nc$ with
$\mathrm{Pl}$ and $\mathrm{Bel}$ throughout.  The five-number scorecard
extends naturally: $\alpha^{*}$ becomes the plausibility
of the observed outcome, and $\eta$ measures the diffuseness of the
plausibility function.  We note
this connection as a direction for future work, as a full treatment
would require addressing the richer combinatorial structure of general
mass functions.

%----------------------------------------------------------
\subsection{Applicability Beyond Severe Weather}%
\label{sec:beyond_severe}

The verification framework is domain-agnostic: any system that
produces subnormal possibility distributions over a finite universe of
discourse $\Omega$ can be evaluated using the five-number scorecard
and the possibility-to-probability bridge
without modification.  The SPC convective outlook categories
(\spc{NONE} through \spc{HIGH}) served as an illustrative running
example, but the same machinery applies directly to air-quality
exceedance categories (e.g., ozone or PM$_{2.5}$ index tiers),
wildfire risk levels, flood inundation severity classes, and climate
scenario classifications.  In each case the forecaster supplies a
possibility distribution $\pi$ over the domain-specific categories,
and the scorecard quantifies how well that distribution captured the
observed outcome.

Extension to continuous forecast variables is straightforward via
$\alpha$-cut discretization.  The continuous possibility distribution
is sliced at a set of threshold levels $\alpha_1 > \alpha_2 > \cdots
> \alpha_k$; each slice defines a nested categorical event (the set of
outcomes with $\pi(\omega) \geq \alpha_i$) for which the scorecard
applies.  For spatial fields, pointwise application
of the scorecard at each grid cell produces spatially resolved
verification maps, while aggregation across grid points yields
domain-averaged skill metrics analogous to those used in conventional
probabilistic forecast evaluation.

%----------------------------------------------------------
\subsection{Generating Possibilistic Forecasts}%
\label{sec:generation}

A natural question arises: where do possibilistic forecasts come from?
This paper addresses verification, not generation, but several pathways
exist for producing the subnormal distributions the framework evaluates.

\textbf{Fuzzy inference systems.}  The most direct route uses fuzzy
rules whose membership functions are possibility distributions by
construction \citep{Klir1995-va}.  A fuzzy inference system maps input
meteorological variables through rule sets with fuzzy antecedents and
consequents; the aggregated output is naturally subnormal when rules
fire at partial strength, encoding the system's incomplete confidence
without normalization.

\textbf{Ensemble-derived possibility.}  For each category, compute the
fraction of ensemble members predicting it---but do \textit{not}
normalize to $\max = 1$.  Let the ensemble spread speak: if members are
scattered across categories, the maximum fraction is low and $\ign$ is
high; if members cluster, the maximum is high and $\ign$ is low.  The
raw ensemble ``support'' for each outcome is already a possibility-like
quantity.  This is the most accessible pathway for a NWP audience:
ensemble spread, formalized as subnormality.

\textbf{Severity--confidence matrices.}  Operational warning frameworks
already encode possibilistic structure.  The UK Met Office National
Severe Weather Warning Service (NSWWS) uses a two-dimensional matrix
of likelihood and impact to determine warning levels
\citep{Neal2014-MISSING}.  The impact axis encodes what \textit{might}
happen (the shape of the possibility distribution), while the likelihood
axis encodes how confident the forecaster is (akin to subnormality).
Figure~\ref{fig:severity_matrix} adapts this concept to the paper's SPC
framework, mapping severity categories against confidence levels to show
how the joint signal varies from ``monitor'' (low severity, high
ignorance) to ``act now'' (high severity, low ignorance).  The
verification framework developed here formalizes this operational
intuition with precise metrics.

\textbf{Conformal prediction.}  Conformal methods produce prediction
sets at given confidence levels---structurally, these are $\alpha$-cuts
of a possibility distribution.  Building sets at multiple $\alpha$
levels constructs a full distribution, and the coverage guarantee maps
onto conditional necessity.

\begin{figure}[t]
\centering
\includegraphics[width=\columnwidth]{severity_matrix.png}
\caption{Severity--confidence matrix adapted from the UK Met Office
    NSWWS framework \citep{Neal2014-MISSING} into the SPC categorical
    system.  Rows are SPC severity categories (\spc{NONE}--\spc{HIGH});
    columns represent forecaster confidence, mapped to subnormality
    ($\ign$).  Color intensity encodes the joint risk signal: deep purple
    indicates high severity with high confidence (act now); pale tones
    indicate low severity or low confidence (monitor).  The matrix
    illustrates how operational warning decisions already implicitly
    navigate the two-dimensional space that possibility theory
    formalizes.}
\label{fig:severity_matrix}
\end{figure}

To connect the matrix to the verification framework, consider a cell at
\spc{MDT} severity with high confidence ($\ign = 0.10$).  The implied
distribution might be
$\pi = (0.00,\, 0.00,\, 0.05,\, 0.15,\, 0.90,\, 0.10)$---exactly
Scenario~A from Section~\ref{sec:worked_examples}.  The five-number
scorecard then evaluates whether that matrix cell's implied distribution
is well-calibrated: $\alpha^* = 1.0$ and $\delta = +0.778$ when
\spc{MDT} verifies, but $\alpha^* = 0$ and $\delta < 0$ when it does
not.  Aggregating scorecard values across all days that fall in a given
matrix cell reveals whether the operational severity--confidence mapping
is empirically justified.

%%% §7.4 (LLM/Ffion communication pipeline) removed in Round 2.
%%% Moved to future-work paragraph below.

%----------------------------------------------------------
\subsection{Software Availability}\label{sec:software}

All verification metrics, the bridge conversion, and the visualization
tools described in this paper are implemented in Python and publicly
available at
\url{https://github.com/bingham-research-center/possibility-verif}.

%----------------------------------------------------------
\subsection{Limitations and Future Directions}%
\label{sec:limitations}

Several limitations should be noted.  Sample-size requirements for the
three reliability hypothesis tests
(Section~\ref{sec:three_tests}) have not been established; at pilot
scales (${\sim}100$ cases), these tests should be supplemented with
illustrative case studies and bootstrap confidence intervals to guard
against small-sample artifacts.  The scorecard is demonstrated here without a
possibilistic skill score, because defining a principled possibilistic
climatology---the baseline needed for skill---remains an open problem.
A consonant (nested) reference distribution derived from observed
category frequencies is one candidate, but it may not suit applications
with multimodal climatologies where the nested-set assumption breaks
down.  Formal credal-set exceedance bounds---grounded in the
possibility--necessity duality but rigorously derived for subnormal
distributions---and construction of possibilistic baselines are
deferred to future work.

The possibility-to-probability bridge assumes that the ignorance
outcome is never observed: the true category is always one of the
elements of $\Omega$.  This is a closed-world assumption, appropriate
when $\Omega$ is exhaustive.  For open-world applications in which the
true outcome may lie outside the defined categories---for instance,
an unprecedented severe weather mode not represented in the SPC
taxonomy---the framework would need modification to allow the ignorance
outcome to verify as ``correct'' under certain conditions, effectively
treating $\ign$ as a legitimate forecast of novelty.

Several directions for future work follow naturally.  The
three-component representation ($\poss$, $\ign$, $\Nc$) lends itself to
downstream communication; a companion project (Lawson, in prep.)
explores LLM-mediated translation of the possibilistic triplet into
tiered plain-language advisories, using the ignorance signal to trigger
hedged language and conditional necessity to trigger confident language.
First,
multivariate possibilistic verification for correlated hazards (e.g.,
simultaneous wind and hail occurrence) would require joint possibility
distributions and an extension of the scorecard to multidimensional
$\Omega$.  Second, conformal prediction also produces interval-valued
forecasts but from a frequentist calibration perspective; a formal
comparison of conformal coverage guarantees with possibilistic
intervals could clarify the relative merits of each
paradigm.  Third, real-time operational implementation with automated
scorecard dashboards would enable forecasters to monitor possibilistic
skill across successive forecast cycles.  Finally, the performance
diagram (Figure~\ref{fig:performance_diagram}) could be extended to
track model development trajectories, plotting successive model
versions as connected points to visualize the trade-offs between
sharpness, calibration, and ignorance over time.
